% Преамбула для дипломного проекта БНТУ
% Соответствует требованиям методички 2026 г. и нормоконтроля

% --- Базовые настройки ---
% XeLaTeX: fontenc и inputenc не нужны — UTF-8 и кириллица поддерживаются нативно
\usepackage[russian]{babel}

% --- Переносы слов ---
% Глобально отключены: слово не разбивается, LaTeX растягивает пробелы.
% В таблицах локально разрешены через \AtBeginEnvironment{longtblr} ниже —
% но фактически tabularray X-ячейки всё равно не вызывают TeX-хифенацию для
% русских слов (известное ограничение). Поэтому в проблемных таблицах
% ширины колонок подобраны под самые длинные неразрывные слова.
\hyphenpenalty=10000
\exhyphenpenalty=10000
\tolerance=9999

% --- Защита от висячих строк (методичка п. 4.1.30) ---
% widowpenalty: одна последняя строка абзаца на следующей странице
% clubpenalty: одна первая строка абзаца внизу страницы
% 10000 = категорически запретить; LaTeX перенесёт лишнюю строку при возможности
\widowpenalty=10000
\clubpenalty=10000

% --- Поля страницы (методичка п. 4.1.1) ---
% левое 30 мм, правое 15 мм, верх/низ 20 мм
\usepackage[a4paper,
            left=30mm, right=15mm,
            top=20mm, bottom=20mm,
            headsep=8mm,
            footskip=10mm]{geometry}

% --- Шрифт Times New Roman 14pt (настоящий системный шрифт, не клон) ---
\usepackage{fontspec}
\setmainfont{Times New Roman}

% Размер шрифта основного текста = 14pt; межстрочный интервал 18pt
\renewcommand{\normalsize}{\fontsize{14pt}{18pt}\selectfont}
\normalsize

% --- Абзацный отступ 1.25 см, без вертикального промежутка между абзацами ---
\usepackage{indentfirst}
\setlength{\parindent}{1.25cm}
\setlength{\parskip}{0pt}

% --- Графика ---
\usepackage[demo]{graphicx}  % demo: заглушки вместо картинок; убери [demo] когда добавишь реальные изображения
\graphicspath{{images/}}

% --- Таблицы ---
\usepackage{array}
\usepackage{longtable}
\usepackage{booktabs}
\usepackage{tabularx}
% tabularray: современный пакет, корректно считает X-колонки за один проход,
% без двухпроходной сходимости через aux (в отличие от ltablex/xltabular)
\usepackage{tabularray}
\UseTblrLibrary{counter}
% Методичка БНТУ:
%   «Таблица 6.1 -- Название» — слева над таблицей, без точки в конце
%   «Продолжение таблицы 6.1» — на следующих страницах, без названия и без «(Continued)»
\DefTblrTemplate{caption-tag}{default}{Таблица~\thetable}
\DefTblrTemplate{caption-sep}{default}{~--~}
\SetTblrTemplate{caption-tag}{default}
\SetTblrTemplate{caption-sep}{default}
% Подпись слева, после неё — 1 свободная строка (18pt) до таблицы (п. 4.1.15)
\DefTblrTemplate{caption}{bntu}{%
  \UseTblrAlign{caption}\UseTblrIndent{caption}\UseTblrHang{caption}%
  \leavevmode%
  \UseTblrTemplate{caption-tag}{default}%
  \UseTblrTemplate{caption-sep}{default}%
  \UseTblrTemplate{caption-text}{default}%
  \par\vspace{18pt}%
}
\SetTblrTemplate{caption}{bntu}
% Для продолжения — отдельный шаблон без названия и без «(Continued)», тоже с отступом
\DefTblrTemplate{capcont-text}{bntu}{Продолжение таблицы~\thetable}
\DefTblrTemplate{capcont}{bntu}{%
  \UseTblrAlign{capcont}\UseTblrIndent{capcont}%
  \UseTblrTemplate{capcont-text}{bntu}\par\vspace{18pt}%
}
\SetTblrTemplate{capcont}{bntu}
% Убрать «Continued on next page» в нижнем колонтитуле многостраничной таблицы
\renewcommand{\tblrcontfootname}{}
% Выравнивание подписи и продолжения — по левому краю
\SetTblrStyle{caption}{halign=l}
\SetTblrStyle{capcont}{halign=l}
% В таблицах:
%   - \code переопределяется локально на seqsplit-вариант, чтобы длинные URL
%     и идентификаторы могли разрываться (вне таблиц \code не разбивается);
%   - \exhyphenpenalty=50 разрешает разрыв по уже существующим дефисам
%     (например, «Бизнес-аналитик» → «Бизнес-» / «аналитик»), но не вводит
%     новых переносов внутри слов;
%   - \hyphenpenalty оставлен на 10000 — авто-перенос русских слов tabularray
%     всё равно не вызывает, см. комментарий в начале преамбулы.
\usepackage{seqsplit}
\AtBeginEnvironment{longtblr}{%
  \exhyphenpenalty=50\tolerance=9999\emergencystretch=3em%
  \let\code\seqsplit%
}
\usepackage{multirow}
\usepackage{makecell}

% --- Списки: тире для маркированных, скобка для нумерованных ---
\usepackage{enumitem}
\setlist[itemize]{label=--, wide, labelsep=0.5em, nosep, topsep=0pt}
\setlist[enumerate]{label=\arabic*), wide, labelsep=0.5em, nosep, topsep=0pt}

% --- Заголовки разделов/подразделов ---
\usepackage{titlesec}

% Раздел: ПРОПИСНЫЕ полужирный 14pt по центру, без точки.
% После заголовка - пустая строка (одна пробельная)
\titleformat{\section}[block]
  {\normalsize\bfseries}
  {\thesection}{0.5em}{\MakeUppercase}
\titlespacing*{\section}
  {1.25cm}{0pt}{18pt}

% Подраздел: строчные с первой прописной, полужирный 13pt (методичка п. 4.1.2)
\titleformat{\subsection}[block]
  {\fontsize{13pt}{16pt}\selectfont\bfseries}
  {\thesubsection}{0.5em}{}
\titlespacing*{\subsection}
  {1.25cm}{18pt}{18pt}

% Пункт (subsubsection): полужирный 13pt — как подраздел.
% Методичка явно задаёт размеры только для разделов (14 пт) и подразделов (13 пт);
% о пунктах не упоминает. Используем 13 пт = размер подраздела, чтобы не вводить
% произвольных значений; иерархия читается по номеру (7.2 → 7.2.1).
\titleformat{\subsubsection}[block]
  {\fontsize{13pt}{16pt}\selectfont\bfseries}
  {\thesubsubsection}{0.5em}{}
\titlespacing*{\subsubsection}
  {1.25cm}{18pt}{18pt}

% Заголовки без нумерации (ВВЕДЕНИЕ, ЗАКЛЮЧЕНИЕ, СПИСОК ...) - по центру
\newcommand{\unnumberedsection}[1]{%
  \clearpage
  {\centering\textbf{\MakeUppercase{#1}}\par}%
  \addcontentsline{toc}{section}{\MakeUppercase{#1}}%
  \vspace{18pt}\noindent\ignorespaces
}

% Каждый раздел начинается с новой страницы
\let\oldsection\section
\renewcommand{\section}{\clearpage\oldsection}

% Заголовки разделов в оглавлении — прописными буквами
\makeatletter
\let\bntu@sec\section
\def\bntu@sec@witharg[#1]#2{\bntu@sec[\MakeUppercase{#1}]{#2}}
\renewcommand{\section}{%
  \@ifstar{\bntu@sec*}%
  {\@dblarg\bntu@sec@witharg}%
}
\makeatother

% --- Подписи к рисункам и таблицам ---
\usepackage{caption}
% Рисунок: «Рисунок 1.1 — Название», по центру, без точки
\DeclareCaptionLabelFormat{figuredash}{Рисунок #2}
\DeclareCaptionLabelSeparator{dash}{~--~}
\captionsetup[figure]{
  labelformat=figuredash,
  labelsep=dash,
  format=plain,
  justification=centering,
  singlelinecheck=false,
  font={normalsize},
  skip=18pt
}
% Отступы вокруг плавающих объектов — 1 свободная строка (п. 4.1.11)
\setlength{\textfloatsep}{18pt}   % float вверху/внизу страницы ↔ текст
\setlength{\intextsep}{18pt}      % float [h] ↔ текст сверху и снизу
\setlength{\floatsep}{18pt}       % между двумя подряд идущими float-ами
% Таблица: «Таблица 1.1 — Название», слева над таблицей
\DeclareCaptionLabelFormat{tabledash}{Таблица #2}
\captionsetup[table]{
  labelformat=tabledash,
  labelsep=dash,
  format=plain,
  justification=raggedright,
  singlelinecheck=false,
  font={normalsize},
  skip=18pt,
  position=top
}
\captionsetup[longtable]{
  labelformat=tabledash,
  labelsep=dash,
  format=plain,
  justification=raggedright,
  singlelinecheck=false,
  font={normalsize},
  skip=18pt,
  position=top
}

% --- Нумерация рисунков, таблиц, формул в пределах раздела ---
\usepackage{chngcntr}
\counterwithin{figure}{section}
\counterwithin{table}{section}
\counterwithin{equation}{section}

% --- Математика ---
\usepackage{amsmath}
\usepackage{amssymb}
\usepackage{amsfonts}

% --- Math-шрифт под Times New Roman ---
% TG Termes Math — Times-клон от GUST для math-режима (математические символы,
% греческие буквы, дробная черта и т.д. оформлены в стиле Times).
% Цифры дополнительно подменяются на реальный Times New Roman, чтобы они
% побайтово совпадали с цифрами в обычном тексте и в таблицах.
\usepackage{unicode-math}
\setmathfont{texgyretermes-math.otf}
\setmathfont{Times New Roman}[range=up/num]

% --- Гиперссылки (для удобства просмотра PDF, не печатается цветом) ---
\PassOptionsToPackage{hyphens}{url}  % перенос длинных URL в библиографии
\usepackage[unicode, hidelinks]{hyperref}
\urlstyle{rm}  % URL в том же шрифте что и основной текст (Times New Roman)

% --- Микротипографика ---
% Выступание пунктуации за поля (protrusion) отключено — требование нормоконтроля.
% Чтобы включить protrusion (типографски красивее): замени строку ниже на \usepackage{microtype}
\usepackage[protrusion=false]{microtype}
\setlength{\emergencystretch}{3em}

% --- Колонтитулы: номер страницы в правом верхнем углу ---
\usepackage{fancyhdr}
\pagestyle{fancy}
\fancyhf{}
\fancyhead[R]{\thepage}
\renewcommand{\headrulewidth}{0pt}
\renewcommand{\footrulewidth}{0pt}

% Первый стиль (для титульного, задания и т.п. - без номера)
\fancypagestyle{empty}{%
  \fancyhf{}
  \renewcommand{\headrulewidth}{0pt}
}

% --- Оглавление: без абзацного отступа, заголовок ОГЛАВЛЕНИЕ ---
% titletoc не используется и конфликтует с tocloft — убран
\usepackage{tocloft}
% Заголовок ОГЛАВЛЕНИЕ: 14pt полужирный по центру (tocloft по умолчанию — \Large)
% babel/russian перебивает \contentsname своим «Содержание» — нужен \addto\captionsrussian
\addto\captionsrussian{\renewcommand{\contentsname}{\MakeUppercase{Оглавление}}}
\renewcommand{\cfttoctitlefont}{\noindent\hfill\normalsize\bfseries}
\renewcommand{\cftaftertoctitle}{\hfill}
\setlength{\cftbeforetoctitleskip}{0pt}
\setlength{\cftaftertoctitleskip}{18pt}
\setlength{\cftbeforesecskip}{2pt}
% Разделы — без отступа; подразделы — небольшое смещение (требование нормоконтроля)
\setlength{\cftsecindent}{0pt}
\setlength{\cftsecnumwidth}{0pt}
\setlength{\cftsubsecindent}{1.5em}
\setlength{\cftsubsecnumwidth}{0pt}
\setlength{\cftsubsubsecindent}{3em}
\setlength{\cftsubsubsecnumwidth}{0pt}
\renewcommand{\cftsecfont}{\normalfont}
\renewcommand{\cftsubsecfont}{\normalfont}
\renewcommand{\cftsubsubsecfont}{\normalfont}
\renewcommand{\cftsecpagefont}{\normalfont}
\renewcommand{\cftsubsecpagefont}{\normalfont}
\renewcommand{\cftsubsubsecpagefont}{\normalfont}
\renewcommand{\cftdotsep}{1.5}
% Точки-заполнители для всех уровней
\renewcommand{\cftsecleader}{\cftdotfill{\cftdotsep}}
\renewcommand{\cftsubsecleader}{\cftdotfill{\cftdotsep}}
\renewcommand{\cftsubsubsecleader}{\cftdotfill{\cftdotsep}}
% Продолжение многострочных заголовков — под цифрой номера (ГОСТ 2.105)
% \numberline по умолчанию кладёт номер в box фиксированной ширины → hanging indent.
% Заменяем на inline: номер + пробел, продолжение выравнивается по \cft*indent.
\makeatletter
\renewcommand{\numberline}[1]{#1\hspace{1em}}
\makeatother

% --- Листинги программного кода ---
\usepackage{listings}
\usepackage{xcolor}
\lstset{
  basicstyle=\small\rmfamily,
  breaklines=true,
  frame=single,
  numbers=left,
  numberstyle=\tiny,
  keywordstyle=\bfseries,
  commentstyle=\itshape,
  showstringspaces=false,
  tabsize=2,
  captionpos=b
}

% --- Умная запятая в формулах: «0,5» вместо «0, 5» ---
\usepackage{icomma}

% --- Сноски: нумерация начинается заново на каждой странице ---
\usepackage{perpage}
\MakePerPage{footnote}

% ---------------------------------------------------------------
% Автоматический подсчёт рисунков, таблиц, формул и источников
% для строки реферата. По образцу шаблона БГУИР.
%
% Использование:
%   \totfig   — итого рисунков (до заключения включительно)
%   \tottab   — итого таблиц
%   \toteq    — итого формул
%   \totref   — итого источников (считает команда \src)
%   \pageref{lastmainpage} — номер последней страницы заключения
% ---------------------------------------------------------------
\usepackage{lastpage}
\usepackage{etoolbox}

\newcounter{totfigures}
\newcounter{tottables}
\newcounter{totreferences}
\newcounter{totequations}

\providecommand\totfig{?}
\providecommand\tottab{?}
\providecommand\totref{?}
\providecommand\toteq{?}

\makeatletter
\AtEndDocument{%
  \addtocounter{totfigures}{\value{figure}}%
  \addtocounter{tottables}{\value{table}}%
  \addtocounter{totequations}{\value{equation}}%
  \immediate\write\@mainaux{%
    \string\gdef\string\totfig{\number\value{totfigures}}%
    \string\gdef\string\tottab{\number\value{tottables}}%
    \string\gdef\string\totref{\number\value{totreferences}}%
    \string\gdef\string\toteq{\number\value{totequations}}%
  }%
}
\makeatother

% Накапливаем счётчики перед каждым \section (который их сбрасывает)
\pretocmd{\section}{%
  \addtocounter{totfigures}{\value{figure}}\setcounter{figure}{0}%
  \addtocounter{tottables}{\value{table}}\setcounter{table}{0}%
  \addtocounter{totequations}{\value{equation}}\setcounter{equation}{0}%
}{}{}

% \src[key] — вместо \item в списке литературы; автоматически считает источники.
% Необязательный аргумент key задаёт метку: \src[banner] → \label{src:banner}.
% В тексте ссылаться как [\ref{src:banner}] — номер обновляется автоматически.
\newcommand{\src}[1][]{\stepcounter{totreferences}\item\ifx\relax#1\relax\else\label{src:#1}\fi}

% --- Перенос слов через дефис: аппаратно\hyphпрограммный ---
\def\hyph{-\penalty0\hskip0pt\relax}

% --- Знак № (устойчив в формулах) ---
\DeclareRobustCommand{\No}{%
  \ifmmode{\nfss@text{\textnumero}}\else\textnumero\fi}

% --- Кириллические перечисления а), б), в) для enumitem ---
% Использование: \begin{enumerate}[label=\asbuk*)]
\makeatletter
\AddEnumerateCounter{\asbuk}{\@asbuk}{щ}
\makeatother

% --- Окружение для расшифровки переменных формулы («где ...») ---
% Методичка п. 4.1.26: первая строка с абзаца со слова «где» без двоеточия;
% по примеру из методички все последующие строки выравниваются с тем же
% абзацным отступом (а не «по margin 0», как можно подумать из формулировки
% «обычное выравнивание»). Использование: \item для каждой переменной.
\usepackage{calc}
\newenvironment{explanationx}{%
  \par
  \def\item{%
    \par\noindent\hspace{1.25cm}\ignorespaces
  }%
}{%
  \par
}

% --- Макросы для названий технологий ---
\newcommand{\csharp}{C\#}
\newcommand{\fsharp}{F\#}
\newcommand{\cpp}{C++}
\newcommand{\dotnet}{Microsoft .NET}
\newcommand{\netcore}{.NET~8}
\newcommand{\java}{Java}
\newcommand{\javascript}{JavaScript}
\newcommand{\typescript}{TypeScript}
\newcommand{\postgresql}{PostgreSQL}
\newcommand{\aspnet}{ASP.NET~Core}
\newcommand{\reactjs}{React}

% --- Цвета для листингов ---
\definecolor{clrKeyword}{rgb}{0.13,0.13,0.8}
\definecolor{clrComment}{rgb}{0,0.50,0}
\definecolor{clrString}{rgb}{0.60,0,0}

% Стиль для C#
\lstdefinestyle{csharpstyle}{
  language=[Sharp]C,
  morekeywords={yield,var,get,set,from,select,partial,where,async,await,
                record,required,init,with,global,file,scoped,nint,nuint},
  breaklines=true,
  columns=fullflexible,
  basicstyle=\small\rmfamily,
  keywordstyle=\bfseries\color{clrKeyword},
  commentstyle=\itshape\color{clrComment},
  stringstyle=\color{clrString},
  showstringspaces=false,
  numbers=left, numberstyle=\tiny\color{gray},
  frame=single, tabsize=2, captionpos=b
}

% Стиль для JavaScript / JSX
\lstdefinestyle{jsxstyle}{
  language=Java,
  morekeywords={const,let,async,await,export,default,import,from,of,
                class,extends,return,function,=>,null,undefined,true,false,
                useState,useEffect,useCallback,useMemo,useRef},
  breaklines=true,
  columns=fullflexible,
  basicstyle=\small\rmfamily,
  keywordstyle=\bfseries\color{clrKeyword},
  commentstyle=\itshape\color{clrComment},
  stringstyle=\color{clrString},
  showstringspaces=false,
  numbers=left, numberstyle=\tiny\color{gray},
  frame=single, tabsize=2, captionpos=b
}

% Стиль для SQL
\lstdefinestyle{sqlstyle}{
  language=SQL,
  breaklines=true,
  columns=fullflexible,
  basicstyle=\small\rmfamily,
  keywordstyle=\bfseries\color{clrKeyword},
  commentstyle=\itshape\color{clrComment},
  stringstyle=\color{clrString},
  showstringspaces=false,
  numbers=left, numberstyle=\tiny\color{gray},
  frame=single, tabsize=2, captionpos=b
}

% --- Команды для удобства ---
\newcommand{\code}[1]{#1}
